\documentclass[11pt,a4paper]{article}
\usepackage[margin=2.5cm]{geometry}
\usepackage{amsmath,amssymb}
\usepackage{booktabs}
\usepackage{xcolor}
\usepackage{hyperref}
\usepackage{enumitem}
\usepackage{parskip}
\usepackage{titlesec}
\usepackage{mdframed}

% Colour scheme
\definecolor{reviewerblue}{RGB}{0,70,127}
\definecolor{responsegreen}{RGB}{0,100,50}
\definecolor{changebrown}{RGB}{139,69,19}
\definecolor{lightgray}{RGB}{245,245,245}

% Reviewer comment box
\newmdenv[
  backgroundcolor=lightgray,
  linecolor=reviewerblue,
  linewidth=1pt,
  leftmargin=0pt,
  rightmargin=0pt,
  innerleftmargin=10pt,
  innerrightmargin=10pt,
  innertopmargin=6pt,
  innerbottommargin=6pt,
  skipabove=8pt,
  skipbelow=4pt
]{reviewbox}

% Section formatting
\titleformat{\section}{\large\bfseries\color{reviewerblue}}{}{0em}{}[\titlerule]
\titleformat{\subsection}{\normalsize\bfseries\color{reviewerblue}}{}{0em}{}

% Macros
\newcommand{\RC}[1]{\begin{reviewbox}\textit{\textbf{Reviewer Comment:} #1}\end{reviewbox}}
\newcommand{\RS}[1]{\noindent\textcolor{responsegreen}{\textbf{Response:}} #1\par\medskip}
\newcommand{\CH}[1]{\noindent\textcolor{changebrown}{\textbf{Change in manuscript:}} #1\par\bigskip}

\begin{document}

% -----------------------------------------------------------------------
%  HEADER
% -----------------------------------------------------------------------
\begin{center}
  {\LARGE\bfseries Second Revision --- Response to Reviewers}\\[1em]
  {\large Manuscript: \textit{Agentic SABRE: An Uncertainty-Aware Neuro-Symbolic
  Multi-Agent Framework for Adaptive Ransomware Detection}}\\[0.5em]
  {\normalsize Expert Systems With Applications --- Second Revision}\\[0.5em]
  {\normalsize \today}
\end{center}

\hrule
\bigskip

We thank the reviewers for their continued engagement with the manuscript. The
second-round comments have prompted four substantive additions to the paper: a
family-holdout generalisation analysis in the Results section, an adaptive
thresholding discussion in the Limitations section, a structural explanation of
the obstacles to standard adversarial attack experiments, and a new comparative
runtime table in the Discussion. All new text is highlighted in \textbf{bold blue}
in the revised manuscript. We also clarify below the architectural distinction
between Agentic SABRE and existing neuro-symbolic frameworks, and the dataset
imbalance and coverage limitations.

\bigskip
\hrule
\bigskip

% -----------------------------------------------------------------------
%  REVIEWER 2 (SECOND ROUND)
% -----------------------------------------------------------------------
\section{Reviewer \#2 --- Second Round}

\subsection*{R2.1 --- Validation on unseen ransomware families or temporally separated datasets}

\RC{Include validation on unseen ransomware families or temporally separated datasets
to demonstrate real-world generalisation capability.}

\RS{We have conducted multi-seed generalisation experiments across four complementary
evaluation splits of RanSMAP using 6 independent random seeds. The results are reported
in the multi-seed generalisation table and the corresponding AUROC figure in the
revised manuscript (Section~IX, new subsection ``Generalisation Across Evaluation
Splits''). Temporal holdout (AUROC $0.960 \pm 0.005$) and hardware holdout
(AUROC $0.939 \pm 0.005$) match random-stratified performance ($0.961 \pm 0.003$),
confirming that Agentic~SABRE does not overfit to specific recording sessions or
hardware configurations. Family holdout reveals a meaningful generalisation gap
(AUROC $0.707 \pm 0.010$), which we acknowledge honestly: when training families are
absent from the test set, the model must rely on general distributional signals rather
than family-specific telemetry patterns. Crucially, the uncertainty mechanism responds
correctly to this degraded regime---samples from unseen families exhibit elevated
$\sigma_{\mathrm{beh}}$ and are routed to \textsc{ESCALATE} rather than
\textsc{AUTO-CONTAIN}, maintaining analyst oversight precisely where the model is
least reliable.}

\CH{New subsection ``Generalisation Across Evaluation Splits'' added in Section~IX
(Results) with a multi-seed results table (four evaluation splits) and an AUROC
distribution figure (6 seeds).}

%----
\subsection*{R2.2 --- Adaptive thresholding for dynamic deployment environments}

\RC{Add a brief adaptive thresholding discussion for dynamic deployment environments.}

\RS{We have added a new paragraph to Section~XIII (Limitations and Future Work)
addressing this point directly. The paragraph explains that the four triage
thresholds ($\tau$, $\kappa$, $\tau_{\text{high}}$, $\kappa_{\text{low}}$) are
currently optimised offline and are not updated during deployment, creating the risk
of miscalibration under concept drift or family emergence. We propose two concrete
adaptive extensions: (1)~a sliding-window recalibration procedure that
re-optimises the objective $J(\tau,\kappa)$ in Eq.~(12) on a rolling buffer of
analyst-labelled triage decisions; and (2)~a Bayesian online update that maintains
a posterior over $(\tau,\kappa)$ conditioned on observed outcomes. Both extensions
are compatible with the modular structure of Agentic~SABRE---the triage policy is
already separated from the neural agents, so threshold recalibration does not require
any model retraining. We note that the existing Limitations section already
identified the modular structure as enabling ``policy-level updates without
retraining''; the new paragraph substantiates this claim with a concrete adaptive
mechanism.}

\CH{New paragraph on adaptive thresholding added at the end of Section~XIII
(Limitations and Future Work).}

%----
\subsection*{R2.3 --- Adversarial attack experiments (FGSM or PGD) or stronger justification}

\RC{Adversarial attack experiments (FGSM or PGD) in addition to theoretical
perturbation analysis, or stronger justification.}

\RS{We have implemented three operationally realistic behavioural evasion attacks on
RanSMAP and report results in the adversarial robustness table and the corresponding per-family figure
in the revised manuscript (new subsection ``Adversarial Robustness Under Behavioural
Evasion'', inserted before Section~X). As noted, standard gradient-based attacks (FGSM,
PGD) require differentiating through the I/O simulator, which is non-differentiable.
Instead we evaluate: (i)~\textbf{slow-encrypt} (reduces read LBA variance and memory
access counts; MITRE ATT\&CK T1486), (ii)~\textbf{memory-mask} (suppresses memory
entropy and GPA variance statistics), and (iii)~\textbf{full-blend} (combines both),
each at 10 intensity levels $\alpha \in [0, 1]$. At maximum intensity ($\alpha=1.0$),
slow-encrypt reduces detection from $0.879$ to $0.817$ (a $7.1\%$ drop) and
memory-mask reduces it to $0.847$. The escalation rate rises correspondingly (from
$0.088$ to $0.102$ under full slow-encrypt), confirming that the triage mechanism
routes adversarially-perturbed samples to human analysts rather than autonomous
containment. Per-family results show WannaCry is most vulnerable (det$=0.809$ at
$\alpha=0.5$), consistent with its elevated baseline uncertainty
($\sigma_{\mathrm{beh}}=0.091$), while Ryuk and Conti remain above $0.950$.}

\CH{New subsection ``Adversarial Robustness Under Behavioural Evasion'' added before
Section~X (Explainability), with an adversarial robustness table (detection/escalation
rates across attack types and intensities) and a per-family detection figure
(slow-encrypt $\alpha=0.5$).}

%----
\subsection*{R2.4 --- New citation in Section 3.1}

\RC{Cite ``Securing Digital Governance: A Deep Learning and Blockchain Framework for
Malware Detection in IoT Networks'' in Section~3.1 (it focuses on intelligent malware
detection using deep learning-based cybersecurity architectures).}

\RS{We have added the requested citation (\texttt{pawar2024securing}) at the end of
Section~3.1 (Machine Learning for Ransomware Detection). The sentence reads: ``Recent
work has also explored deep learning and blockchain-based frameworks for malware
detection in IoT environments~\cite{pawar2024securing}, underscoring the breadth of
architectural approaches being applied to adaptive threat analysis and explainable
malware defence.'' The full bibliographic entry has been added to the reference list.}

\CH{Citation \texttt{pawar2024securing} added at the end of Section~3.1 with
accompanying sentence. Bibliographic entry added to \texttt{Reference\_Sabre.bib}.}

%----
\subsection*{R2.5 --- Comparative runtime or resource-consumption evaluation}

\RC{A small comparative runtime or resource-consumption evaluation against existing
ransomware detection approaches.}

\RS{We have replaced the theoretical Big-O complexity table with empirically measured
inference latency on an NVIDIA GB10 GPU (500 repetitions, 50 warm-up), reported in
the runtime comparison table in the revised manuscript (Section~XII,
``Measured runtime performance''). Key results: the Behavioral CNN requires
$1.94$\,ms (p50) for $T=10$ MC Dropout passes; the Semantic CNN requires $1.77$\,ms;
and the full pipeline is estimated at $\approx 3.7$\,ms ($\approx 270$ samples/s).
The MC Dropout overhead is $3.7\times$ relative to a single deterministic forward
pass ($0.52$\,ms). The Behavioral CNN has only 8,385 parameters (32.8\,KB), making
it suitable for endpoint deployment. At a 10-second telemetry window, the pipeline
supports over 2,000 concurrent monitored streams before becoming the bottleneck.}

\CH{Theoretical Big-O runtime table replaced with empirically measured latency table
(the runtime comparison table) in Section~XII (Discussion). Paragraph
retitled ``Measured runtime performance'' with GPU benchmark results on NVIDIA GB10.}

\bigskip
\hrule
\bigskip

% -----------------------------------------------------------------------
%  REVIEWER 5 (SECOND ROUND)
% -----------------------------------------------------------------------
\section{Reviewer \#5 --- Second Round}

\subsection*{R5.1 --- Architecture differentiation from existing neuro-symbolic frameworks and necessity of multi-agent design}

\RC{The manuscript does not clearly articulate: (a)~how the architecture differs from
existing neuro-symbolic cybersecurity frameworks, (b)~why the multi-agent design is
necessary rather than architectural overhead.}

\RS{We have added a new paragraph in Section~3.5 (Multi-Agent and Neuro-Symbolic
Approaches) addressing both sub-points directly, inserted immediately before the
existing ``However, these prior systems typically lack\ldots'' paragraph.

\textbf{Differentiation from existing neuro-symbolic systems (R5.1a).}
Prior neuro-symbolic cybersecurity frameworks---including the knowledge-base
integration approaches surveyed by Yang et al.~\cite{YangSBK25} and the neural-symbolic
architectures reviewed by Besold et al.~\cite{Besold2021neurosymbolicsurvey}---embed
symbolic rules at the architecture or knowledge-base level, structurally coupling the
neural and symbolic components so that neither can be updated independently. In
Agentic~SABRE, the symbolic component consists of ATT\&CK-inspired interpretable
triage thresholds (API call category flags, entropy-threshold signals, filesystem
access pattern heuristics) that are evaluated as conditions on the fused risk score
and epistemic uncertainty. These thresholds are encoded as the parameters
$(\tau, \kappa, \tau_{\text{high}}, \kappa_{\text{low}})$ and can be recalibrated
without retraining the neural agents whenever threat intelligence or operational
requirements change. This separation between the learning layer and the symbolic
reasoning layer is a deliberate architectural choice.

\textbf{Necessity of multi-agent design (R5.1b).}
The multi-agent design is motivated by complementary modality specialisation, not
redundancy. The Semantic Agent operates on PE-level embedding vectors
($\mathbb{R}^{384}$) derived from static file metadata, while the Behavioural Agent
operates on sliding-window I/O telemetry vectors ($\mathbb{R}^{16}$) from runtime
memory and storage access. These are structurally heterogeneous input spaces: a
single monolithic agent would either require a common representation that discards
modality-specific structure, or would need to be trained on concatenated features
that conflate semantically incompatible signals. By maintaining separate agents,
SABRE allows each modality to contribute its own posterior probability
\emph{and its own epistemic uncertainty}, which are then combined by the
orchestrator. The uncertainty from the weaker behavioural modality is the primary
mechanism by which the system escalates ambiguous cases---a function that cannot
be replicated by a single-agent architecture. The ``What Makes This Different from
Ensemble-Based Approaches'' subsection in Section~I further elaborates this point
as addressed in the first-round revision.}

\CH{New differentiating paragraph added in Section~3.5, immediately before the
existing ``However\ldots'' paragraph. The ``What Makes This Different from
Ensemble-Based Approaches'' subsection in Section~I remains as added in the
first revision.}

%----
\subsection*{R5.2 --- Dataset issues: description, modern threat coverage, imbalance, and contamination}

\RC{Dataset Issues: datasets not sufficiently described; no evidence ransomware
samples represent modern evolving threats; no discussion of dataset imbalance and
contamination risks.}

\RS{We have added a new paragraph titled \textbf{``Dataset characteristics,
imbalance, and contamination''} in Section~VIII-A (Datasets), immediately after
the existing ``CTGAN labels'' paragraph, addressing all three sub-issues.

\textbf{Dataset description and imbalance.}
RDset is described with its full sample counts: 138,047 PE samples (41,123
ransomware, 96,924 benign), with $\approx$30\% ransomware prevalence constituting
heavy class imbalance. RanSMAP is described with its 11,820 execution traces from
24 ransomware families and benign workloads across diverse hardware, with imbalance
addressed via SMOTE at the feature level.

\textbf{Contamination risk.}
For RDset, training and test splits are disjoint at the binary (file hash) level:
no sample-level contamination is present, though the absence of a family-wise split
means structural similarity bias is possible. For RanSMAP, contamination risk is
low: splits are defined by disjoint session identifiers, ensuring no execution
session appears in both training and test sets.

\textbf{Modern threat coverage.}
We explicitly acknowledge that neither dataset includes post-2021 ransomware families
such as Cl0p, BlackCat/ALPHV, or LockBit 3.0. This is identified as a known
coverage limitation, and generalisation to modern families is stated as a priority
direction for future work. This is consistent with the temporal coverage constraints
noted in the first-round revision (Reviewer~\#2, Point~2) and the ``Generalisation
caveats'' paragraph in Section~IX.}

\CH{New \textbf{``Dataset characteristics, imbalance, and contamination''} paragraph
added in Section~VIII-A (Datasets), immediately after the ``CTGAN labels'' paragraph.}

\bigskip
\hrule
\bigskip

% -----------------------------------------------------------------------
%  CLOSING
% -----------------------------------------------------------------------
\section*{Closing Remark}

We believe the second revision addresses all remaining reviewer concerns in a precise
and targeted manner. The additions are surgical: each new paragraph is confined to the
specific location where the relevant evidence or argument belongs, without disturbing
the structure of the existing manuscript. We are grateful to the reviewers for the
rigour of their second-round evaluation and remain available to address any further
questions the editor or reviewers may have.

\bigskip
\noindent\textit{The Authors}

\bibliographystyle{unsrt}
\bibliography{Reference_Sabre}

\end{document}
